\documentclass[10pt,a4paper]{article}
\usepackage[utf8]{inputenc}
\usepackage[T1]{fontenc}
\usepackage{lmodern}
\usepackage[margin=2.3cm]{geometry}
\usepackage{microtype}
\usepackage{amsmath,amssymb}
\usepackage{booktabs,longtable,array}
\usepackage{graphicx}
\usepackage[font=small,labelfont=bf,skip=4pt]{caption}
\usepackage{xcolor}
\usepackage[hidelinks]{hyperref}
\usepackage[section]{placeins}
\usepackage[round]{natbib}
\usepackage{adjustbox}

\graphicspath{{figures/}}
\newcommand{\DataDate}{2026-10-14}
\newcommand{\FreezeDate}{2026-10-14}
\newcommand{\FirstMonth}{2006-02}
\newcommand{\LastMonth}{2026-09}
\newcommand{\NMonths}{248}
\newcommand{\NFormations}{249}
\newcommand{\AvgUniverse}{31.6}
\newcommand{\NTickers}{58}
\newcommand{\NMembersEver}{65}
\newcommand{\NMissing}{7}
\newcommand{\CoverageMean}{96.9\%}
\newcommand{\NChanges}{43}
\newcommand{\CostBps}{10}
\newcommand{\BorrowBps}{25}
\newcommand{\NTrials}{36}
\newcommand{\SRStar}{0.43}
\newcommand{\DaxCagr}{7.4\%}
\newcommand{\DaxSharpe}{0.43}
\newcommand{\DaxMaxDD}{$-$52.4\%}
\newcommand{\SpxCagr}{11.5\%}
\newcommand{\SpxSharpe}{0.74}
\newcommand{\CompCagr}{9.8\%}
\newcommand{\CompSharpe}{0.67}
\newcommand{\SurvBias}{3.8\%}
\newcommand{\NBigMoves}{23}
\newcommand{\NFixes}{9}
\newcommand{\FindBest}{The strongest signal is \emph{Momentum 12-1}, whose long-short portfolio earns 7.2\% per year after costs (Sharpe 0.40); its mean return is only marginally significant (10\% level).}
\newcommand{\FindWorst}{The weakest is \emph{Short-term reversal} with a net Sharpe ratio of $-$0.29; it loses money mostly through trading costs, since its gross spread is only $-$0.4\% per year at a turnover of 16.3 times the portfolio per year.}
\newcommand{\FindHolm}{After a Holm correction for testing six strategies, no long-short premium remains significant at the 5\% level.}
\newcommand{\FindDSR}{Across all 36 specifications that were run, an unskilled researcher would expect a best Sharpe ratio of 0.43 by luck alone. No strategy has a deflated Sharpe ratio above 0.95, so none of them can be called a discovery on this data.}
\newcommand{\FindBeta}{Five of six long-short portfolios carry a negative market beta (lowest: Low beta at $-$0.76), because the bottom tercile of these signals is dominated by volatile, high-beta stocks. In a rising market that drag hides part of the premium: after controlling for the DAX, the CAPM alpha is significant at 5\% for Momentum 12-1, Low beta and 52-week high.}
\newcommand{\FindComposite}{The long-only composite portfolio returns 9.8\% per year against 7.4\% for the DAX, with a beta of 0.65, a tracking error of 9.7\% and an information ratio of 0.16. Its CAPM alpha of 4.2\% per year has a $t$-statistic of 2.80.}
\newcommand{\FindSurvivorship}{Using today's members for the whole history raises the equal-weight return from 7.1\% to 11.0\% per year, a bias of 3.8\% per year that would be mistaken for skill.}
\newcommand{\FindSPX}{Measured in euros, the S\&P 500 compounded at 11.5\% per year versus 7.4\% for the DAX, with lower volatility (14.4\% vs 17.7\%). The correlation of monthly returns is 0.69 and the S\&P 500's beta to the DAX is 0.56, so US equities diversified a euro investor's DAX risk rather than replicating it.}
\newcommand{\FindDefensive}{The defensive signals (low volatility and low beta) give long-only portfolios with an average volatility of 13.0\% against 17.7\% for the DAX and an average Sharpe ratio of 0.64, while their long-short spreads average a Sharpe ratio of 0.03: in raw returns their benefit is lower risk rather than a dollar-neutral premium.}


\newcommand{\fig}[3][\linewidth]{%
  \begin{figure}[htbp]\centering\includegraphics[width=#1]{#2}\caption{#3}\label{fig:#2}\end{figure}}
\newcommand{\tabfile}[3]{%
  \begin{table}[htbp]\centering\caption{#2}\label{tab:#1}\small\begin{adjustbox}{max width=\linewidth}\input{generated/#1.tex}\end{adjustbox}\par\smallskip
  \begin{minipage}{\linewidth}\footnotesize #3\end{minipage}\end{table}}

\title{\textbf{Factor Premia in the DAX}\\[2pt]\large A point-in-time study of German large caps, benchmarked against the S\&P\,500}
\author{Tristan Hartenstein}
\date{Data through \DataDate{} \textperiodcentered{} Specification frozen \FreezeDate}

\begin{document}
\maketitle

\begin{abstract}\noindent
I test five well-documented price-based equity factors (momentum, short-term reversal, low volatility,
low beta and the 52-week high) and an equal-weight composite inside the DAX, from \FirstMonth{} to
\LastMonth{} (\NMonths{} months). The universe is rebuilt day by day from the \NChanges{} documented index
changes, so every portfolio only holds stocks that were in the index at the time.
\FindBest{} \FindHolm{} \FindDSR{} \FindComposite{} \FindSurvivorship{}
All numbers in this paper are produced by code and update automatically when new data arrives.
\end{abstract}

\section{Motivation}
Factor studies on the DAX are easy to get wrong. The index has only 30 members until September 2021 and 40
afterwards, so a few stocks drive every result. Free data contains only today's survivors, which turns
the past into a story of winners. Finally, with enough variants, something always ``works''. This study
addresses each problem directly:
\begin{enumerate}
  \item \textbf{Point-in-time universe.} Membership is reconstructed from the current constituent list and
  every inclusion and exclusion since 2005. Coverage gaps from delisted companies are measured and
  reported, not hidden (Section~\ref{sec:data}).
  \item \textbf{Pre-specified design.} All factor definitions follow the original papers. Parameters are
  fixed in one configuration file and were not tuned; the specification was frozen on \FreezeDate.
  \item \textbf{Honest inference.} Newey--West $t$-statistics, the probabilistic and deflated Sharpe
  ratio \citep{bailey2014}, Holm-adjusted $p$-values and a robustness grid of \NTrials{} variants that all
  count as trials.
  \item \textbf{Realistic implementation.} Trades happen one day after the signal, cost
  \CostBps{}~bps per side, and short positions pay \BorrowBps{}~bps p.a. in lending fees.
  \item \textbf{Live record.} Everything after the freeze date is genuine out-of-sample data
  (Section~\ref{sec:live}).
\end{enumerate}

\section{Data and universe}\label{sec:data}
\paragraph{Sources.} Daily total-return prices (adjusted for dividends and splits) come from Yahoo
Finance for \NTickers{} stocks. The benchmark is the DAX performance index, which reinvests dividends.
The S\&P\,500 total-return index is converted into euros at the daily EUR/USD rate. The risk-free rate is
the German 3-month interbank rate from FRED. Index changes are taken from the Wikipedia list of DAX
changes and checked mechanically: the reconstructed index must contain exactly 30 members on every
day before 20 September 2021 and 40 afterwards.

\paragraph{Coverage.} \NMembersEver{} companies were DAX members at some point since 2005. For \NMissing{}
of them no free price history exists (Table~\ref{tab:tab_missing}), mostly takeovers and insolvencies
before 2013. On average \CoverageMean{} of the index is covered (Figure~\ref{fig:coverage}). The missing
names are excluded from the universe while they were members. This is the remaining survivorship
bias, and it is small but not zero: Wirecard's collapse in 2020, for example, is not in the data.

\fig{coverage}{DAX members and members with usable price data. The shaded area is the gap left by
delisted companies without free data.}
\tabfile{tab_missing}{Index members without free price data}{These stocks are excluded while they were
members, so the effective universe is slightly smaller than the index in those periods.}

\paragraph{Cleaning.} Free data has defects. Each correction is applied explicitly and logged
(Table~\ref{tab:tab_quality}). One share split was missing from the Yahoo series. On spin-off ex-dates
Yahoo's adjustment for German stocks is inconsistent, so the parent's return on that day is replaced by
the DAX return; this removes the error at the cost of one day of stock-specific information. Jumps of
more than 20\% that reverse the next day are treated as bad ticks. All other large moves are kept:
\NBigMoves{} single-day moves above 25\% remain in the sample after inspection, such as the Volkswagen
short squeeze of October 2008 or the Siemens Energy profit warnings of 2023.

\tabfile{tab_quality}{Data-quality corrections}{Raw = return in the Yahoo series, Clean = return
used in the study.}

\paragraph{Survivorship bias.} To show why the point-in-time universe matters, I compare an equal-weight
portfolio of the actual members with one that simply uses today's members for the whole history
(Figure~\ref{fig:survivorship}). \FindSurvivorship{}

\fig{survivorship}{Equal-weight portfolios, monthly rebalanced, after costs. The grey line uses today's
40 members back-filled to 2006 and is therefore biased.}

\section{The market: DAX versus S\&P\,500}
Before testing factors it helps to know the market they live in. \FindSPX{} Table~\ref{tab:tab_benchmarks}
reports the full set of risk and return measures, and Figures~\ref{fig:index_growth} to
\ref{fig:dax_distribution} show the time series, rolling risk, calendar returns, the month-by-month
history and the fat-tailed distribution of daily DAX returns.

\tabfile{tab_benchmarks}{Benchmarks, \FirstMonth{} to \LastMonth{}}{Monthly data. Sharpe ratios use the
German 3-month rate; the standard error follows \citet{mertens2002}, which allows for skewness and fat
tails. VaR and expected shortfall are historical, one-month, 95\%. ``Equal-weight PIT'' holds all
point-in-time members with equal weights and pays the same trading costs as the strategies.}

\fig{index_growth}{Total return in euros and drawdowns. Legend values are the final level of 100
invested at the start.}
\fig{rolling_risk}{Rolling one-year volatility from daily returns and rolling three-year correlation of
monthly returns.}
\fig{calendar_returns}{Calendar-year returns.}
\fig{dax_monthly_heatmap}{Every monthly DAX return of the sample.}
\fig{dax_distribution}{Daily DAX returns against a normal distribution with the same mean and
volatility. VaR and expected shortfall at 99\% are historical.}

\section{Methodology}
\paragraph{Signals.} Each signal is computed at the close of the last trading day of month $t$ from daily
total returns $r_{i,d}$ of the previous $L=252$ trading days. A stock needs at least 200 valid returns in
that window and must be a DAX member on the formation day. Higher values mean higher expected returns.
\begin{align*}
\text{Momentum 12-1} &= \textstyle\prod_{d=t-251}^{t-21}(1+r_{i,d})-1 &&\text{\citep{jegadeesh1993}}\\
\text{Short-term reversal} &= -\Big(\textstyle\prod_{d=t-20}^{t}(1+r_{i,d})-1\Big) &&\text{\citep{jegadeesh1990}}\\
\text{Low volatility} &= -\sigma_i\sqrt{252} &&\text{\citep{ang2006,blitz2007}}\\
\text{Low beta} &= -\big(0.6\,\hat\beta_i + 0.4\big) &&\text{\citep{frazzini2014}}\\
\text{52-week high} &= P_{i,t}\,/\,\max_{d\in(t-252,t]} P_{i,d} &&\text{\citep{george2004}}
\end{align*}
Here $\hat\beta_i=\operatorname{cov}(r_i,r_{\text{DAX}})/\operatorname{var}(r_{\text{DAX}})$ over the same window, shrunk towards 1. The composite is the average of the five cross-sectional $z$-scores, each winsorised at $\pm3$. Its
weights are equal by construction and were not optimised.

\paragraph{Portfolios.} Stocks are sorted into terciles each month (about 10 per leg with 30 members, 13
with 40). Portfolios are equal-weighted, bought at the close one trading day after formation and held
until the next rebalance, with weights drifting in between. The long-short portfolio is long the top
and short the bottom tercile, so it is self-financing and its return is already an excess return.
The long-only portfolio holds the top tercile and is compared with the DAX.

\paragraph{Costs.} Turnover is measured against drifted weights. Every euro traded costs \CostBps{}~bps,
a realistic all-in figure (spread plus fees) for DAX stocks; the short leg pays \BorrowBps{}~bps per year.
The break-even cost is the one-way cost that would reduce the mean excess return to zero.

\paragraph{Inference.} Means and CAPM alphas are tested with Newey--West standard errors. The rank
information coefficient (IC) is the Spearman correlation between the signal and the next holding-period
return. The probabilistic Sharpe ratio (PSR) is the probability that the true Sharpe ratio exceeds zero
given the sample length, skewness and kurtosis. The deflated Sharpe ratio (DSR) replaces zero with the
Sharpe ratio one would expect from the best of \NTrials{} unskilled trials, here \SRStar{}.

\section{Results}
\FindBest{} \FindWorst{} \FindHolm{}

\tabfile{tab_factor_tests}{Do the signals predict returns?}{IC = mean monthly rank IC, ICIR =
annualised IC/standard deviation. Spread = top minus bottom tercile before costs, per year. L/S net = mean
long-short return after costs and lending fees, per year; stars mark Newey--West significance at 10, 5
and 1\%. Holm $p$ adjusts for six tests. PSR and DSR as defined above. Monot. = average tercile returns
decline from top to bottom.}

\fig{tercile_returns}{Average annual return of each tercile. A useful signal shows a staircase from left
(top) to right (bottom).}
\fig{factor_long_short}{Long-short portfolios, growth of 100 before and after costs.}

\paragraph{Market exposure.} \FindBeta{} Table~\ref{tab:tab_long_short} shows the trade-off between gross
premium, turnover and cost.

\tabfile{tab_long_short}{Long-short portfolios after costs}{Beta and alpha from a regression of monthly
returns on DAX excess returns, alpha per year with Newey--West $t$-statistic. Turnover = one-way
turnover per year (1.0 means the whole portfolio is replaced once). Break-even in bps per side.}

\fig{rolling_ic}{Rolling two-year average of the monthly rank IC. Long stretches above or below zero show
that factor premia come and go.}
\fig[0.72\linewidth]{factor_correlations}{Correlations of monthly long-short returns and the DAX. Highly
correlated signals add little diversification when they are combined in the composite.}

\paragraph{Long only.} Most investors cannot short. Table~\ref{tab:tab_long_only} evaluates the top
tercile alone against the DAX. \FindComposite{}

\tabfile{tab_long_only}{Long-only top-tercile portfolios against the DAX}{After costs. TE = tracking
error, IR = information ratio against the DAX, Up/Down = up- and down-market capture ratios. The
equal-weight PIT portfolio shows how much of any outperformance comes from equal weighting alone.}

\fig{composite_long_only}{Long-only composite against both benchmarks.}

\section{Robustness}
Figure~\ref{fig:robustness} repeats every long-short strategy with no, one and five days between signal
and trade and with terciles or quintiles. All \NTrials{} cells enter the deflated Sharpe ratio, so trying
more variants makes the bar for significance higher, not lower. Figure~\ref{fig:cost_sensitivity} shows
how quickly high-turnover signals die when costs rise, and Table~\ref{tab:tab_subperiods} splits the
sample into three regimes.

\fig{robustness}{Net Sharpe ratio for each specification. The base case is lag 1, terciles.}
\fig{cost_sensitivity}{Sharpe ratio as a function of the one-way trading cost.}
\tabfile{tab_subperiods}{Subperiods: Sharpe ratio / CAGR}{Long-short Sharpe ratios are self-financing;
long-only and benchmark Sharpe ratios are in excess of cash.}

\section{Live out-of-sample record}\label{sec:live}
The specification in \texttt{config.py} was frozen on \FreezeDate. Every observation after that date was
unknown when the study was designed, so this table is the only clean test. It grows with each data
update.

\begin{table}[htbp]\centering\small\emph{No out-of-sample data yet. The specification was frozen on 2026-10-14; this table fills in automatically with every update.}
\end{table}

\section{Conclusion}
Inside the DAX, simple price-based factors behave much like the literature predicts in direction, but the
cross-section is narrow and the evidence is weak once costs, market exposure and multiple testing are
taken seriously. \FindHolm{} \FindDSR{} \FindDefensive{} Two lessons reach beyond factor premia:
how much a back-filled universe flatters every backtest, and how different the DAX and the S\&P\,500
have been for a euro investor. The live record will show whether any of the in-sample patterns
persist.

\paragraph{Limitations.} Free data only; \NMissing{} delisted members are missing; prices are adjusted closes,
so dividends enter the 52-week high; costs are a single flat estimate; no fundamental factors (value,
quality), because point-in-time fundamentals are not freely available.

\appendix
\section{Calendar-year returns}
\begin{table}[htbp]\centering\small\begin{tabular}{lrrrrr}
\toprule
Year & DAX & S\&P 500 (EUR) & Equal-weight & Composite (long) & Momentum (L/S) \\
\midrule
2006 & 15.2\% & 2.9\% & 18.8\% & 18.2\% & 0.7\% \\
2007 & 22.3\% & $-$4.7\% & 20.1\% & 25.7\% & 13.1\% \\
2008 & $-$40.4\% & $-$35.8\% & $-$41.6\% & $-$32.0\% & 30.9\% \\
2009 & 23.8\% & 27.7\% & 29.8\% & 27.7\% & $-$24.6\% \\
2010 & 16.1\% & 23.6\% & 20.4\% & 11.0\% & 26.3\% \\
2011 & $-$14.7\% & 4.2\% & $-$14.4\% & 5.8\% & 7.1\% \\
2012 & 29.1\% & 11.5\% & 28.6\% & 28.6\% & $-$1.8\% \\
2013 & 25.5\% & 29.2\% & 21.6\% & 27.2\% & 19.1\% \\
2014 & 2.7\% & 30.5\% & 3.9\% & 8.6\% & $-$8.1\% \\
2015 & 9.6\% & 12.7\% & 11.3\% & 24.4\% & 15.2\% \\
2016 & 6.9\% & 14.6\% & 7.2\% & 9.0\% & $-$11.7\% \\
2017 & 12.5\% & 7.9\% & 19.0\% & 15.8\% & $-$0.2\% \\
2018 & $-$18.3\% & $-$1.0\% & $-$19.0\% & $-$8.7\% & 8.1\% \\
2019 & 25.5\% & 35.2\% & 22.3\% & 27.8\% & 11.5\% \\
2020 & 3.5\% & 7.7\% & 7.5\% & $-$4.3\% & $-$22.0\% \\
2021 & 15.8\% & 40.1\% & 11.0\% & 12.0\% & $-$2.6\% \\
2022 & $-$12.3\% & $-$13.0\% & $-$16.2\% & $-$13.1\% & 17.8\% \\
2023 & 20.3\% & 21.6\% & 16.8\% & 17.0\% & $-$11.2\% \\
2024 & 18.8\% & 33.2\% & 13.9\% & 18.1\% & 33.8\% \\
2025 & 23.0\% & 4.8\% & 20.5\% & 6.3\% & 47.0\% \\
2026 & 2.9\% & 16.2\% & 2.0\% & 2.0\% & 2.0\% \\
\bottomrule
\end{tabular}
\end{table}

\section{Every constituent}
\fig{stocks_risk_return}{Annualised return and volatility of each stock while it was a DAX member,
\FirstMonth{} onwards. Stocks with less than one year of membership are left out of the chart but are in
the table.}
{\small Statistics only cover the days a stock was in the index. $^{\dagger}$ = no longer a member.\par}
{\footnotesize\begin{longtable}{llrrrrrrr}
\toprule
Company & Sector & Years & Return & Vol. & Sharpe & Beta & Max DD & Worst day \\
\midrule
\endhead
\bottomrule
\endfoot
Rheinmetall & Industrials & 3.6 & 50.2\% & 41.2\% & 1.12 & 1.02 & $-$53.0\% & $-$18.7\% \\
Siemens Energy & Industrials & 5.1 & 49.8\% & 51.8\% & 1.02 & 1.48 & $-$71.9\% & $-$37.3\% \\
Hannover Re & Financials & 5.1 & 23.6\% & 24.0\% & 0.91 & 0.65 & $-$17.5\% & $-$7.7\% \\
Daimler Truck & Industrials & 4.6 & 16.2\% & 30.6\% & 0.56 & 1.14 & $-$34.6\% & $-$8.7\% \\
GEA Group & Industrials & 1.1 & 14.9\% & 22.1\% & 0.64 & 0.68 & $-$16.7\% & $-$5.0\% \\
Linde$^{\dagger}$ & Materials & 17.2 & 14.4\% & 25.1\% & 0.63 & 0.44 & $-$38.2\% & $-$9.5\% \\
Volkswagen (ord.)$^{\dagger}$ & Consumer & 3.9 & 13.0\% & 89.0\% & 0.40 & 0.45 & $-$85.7\% & $-$22.7\% \\
Deutsche Wohnen$^{\dagger}$ & Real Estate & 1.4 & 12.7\% & 26.7\% & 0.60 & 0.33 & $-$17.5\% & $-$8.2\% \\
Munich Re & Financials & 20.8 & 12.6\% & 24.8\% & 0.56 & 0.84 & $-$48.4\% & $-$17.7\% \\
Airbus & Industrials & 5.1 & 12.5\% & 28.8\% & 0.48 & 1.22 & $-$27.7\% & $-$11.6\% \\
Deutsche Börse & Financials & 20.8 & 12.0\% & 29.5\% & 0.49 & 0.84 & $-$76.9\% & $-$12.1\% \\
Allianz & Financials & 20.8 & 10.6\% & 28.9\% & 0.45 & 1.10 & $-$72.7\% & $-$15.3\% \\
Fresenius & Health Care & 17.7 & 10.6\% & 25.6\% & 0.49 & 0.68 & $-$72.3\% & $-$17.7\% \\
Siemens & Industrials & 20.8 & 10.2\% & 29.7\% & 0.43 & 1.15 & $-$67.3\% & $-$17.1\% \\
Salzgitter$^{\dagger}$ & Materials & 1.5 & 9.5\% & 49.4\% & 0.41 & 1.50 & $-$39.9\% & $-$16.1\% \\
SAP & Technology & 20.8 & 9.5\% & 26.0\% & 0.43 & 0.81 & $-$53.0\% & $-$21.9\% \\
Heidelberg Materials & Materials & 16.4 & 9.3\% & 31.2\% & 0.42 & 1.18 & $-$66.6\% & $-$14.2\% \\
DHL Group & Industrials & 20.8 & 8.6\% & 28.0\% & 0.39 & 0.94 & $-$72.1\% & $-$17.3\% \\
Deutsche Telekom & Telecom & 20.8 & 8.5\% & 23.2\% & 0.42 & 0.69 & $-$40.6\% & $-$13.0\% \\
Adidas & Consumer & 20.8 & 7.7\% & 31.0\% & 0.35 & 0.90 & $-$71.5\% & $-$15.4\% \\
BASF & Materials & 20.8 & 7.2\% & 28.6\% & 0.34 & 1.08 & $-$60.3\% & $-$13.6\% \\
MTU Aero Engines & Industrials & 7.1 & 6.8\% & 38.5\% & 0.33 & 1.22 & $-$62.9\% & $-$15.6\% \\
Merck KGaA & Health Care & 19.5 & 6.6\% & 25.9\% & 0.34 & 0.63 & $-$52.6\% & $-$14.7\% \\
Henkel & Consumer & 20.8 & 6.6\% & 22.6\% & 0.34 & 0.62 & $-$56.3\% & $-$11.3\% \\
BMW & Consumer & 20.8 & 5.9\% & 30.9\% & 0.30 & 1.11 & $-$65.4\% & $-$13.2\% \\
Beiersdorf & Consumer & 17.1 & 5.2\% & 20.7\% & 0.31 & 0.49 & $-$53.3\% & $-$20.1\% \\
Mercedes-Benz Group & Consumer & 20.8 & 5.2\% & 33.5\% & 0.28 & 1.27 & $-$76.5\% & $-$18.9\% \\
Volkswagen (pref.) & Consumer & 16.9 & 4.7\% & 33.3\% & 0.28 & 1.23 & $-$63.0\% & $-$19.8\% \\
Bayer & Health Care & 20.8 & 4.1\% & 30.4\% & 0.24 & 0.93 & $-$82.3\% & $-$18.0\% \\
RWE & Utilities & 20.8 & 3.4\% & 30.3\% & 0.22 & 0.81 & $-$85.4\% & $-$17.3\% \\
E.ON & Utilities & 20.8 & 3.1\% & 27.7\% & 0.20 & 0.82 & $-$77.3\% & $-$11.5\% \\
Metro (old)$^{\dagger}$ & Consumer & 6.7 & 2.9\% & 36.7\% & 0.20 & 0.80 & $-$71.5\% & $-$13.9\% \\
Infineon & Technology & 20.3 & 2.6\% & 43.3\% & 0.25 & 1.26 & $-$97.1\% & $-$39.6\% \\
Fresenius Medical Care & Health Care & 19.1 & 2.3\% & 24.9\% & 0.18 & 0.49 & $-$69.1\% & $-$16.5\% \\
Lufthansa$^{\dagger}$ & Industrials & 14.5 & 1.4\% & 33.9\% & 0.18 & 0.97 & $-$75.2\% & $-$14.2\% \\
Vonovia & Real Estate & 11.1 & 0.5\% & 28.8\% & 0.13 & 0.65 & $-$65.8\% & $-$12.1\% \\
Continental & Consumer & 17.0 & 0.4\% & 35.0\% & 0.15 & 1.12 & $-$76.7\% & $-$17.5\% \\
thyssenkrupp$^{\dagger}$ & Materials & 13.7 & $-$1.4\% & 39.1\% & 0.13 & 1.31 & $-$74.7\% & $-$12.5\% \\
Qiagen & Health Care & 5.1 & $-$1.6\% & 24.1\% & $-$0.04 & 0.42 & $-$43.4\% & $-$10.8\% \\
Deutsche Bank & Financials & 20.8 & $-$1.6\% & 41.0\% & 0.13 & 1.43 & $-$92.7\% & $-$18.4\% \\
Brenntag & Materials & 5.1 & $-$2.1\% & 26.5\% & $-$0.03 & 0.82 & $-$44.7\% & $-$9.6\% \\
Covestro$^{\dagger}$ & Materials & 6.8 & $-$2.1\% & 34.1\% & 0.08 & 1.04 & $-$69.1\% & $-$14.6\% \\
Symrise & Materials & 5.1 & $-$2.4\% & 24.0\% & $-$0.07 & 0.47 & $-$47.0\% & $-$9.1\% \\
Siemens Healthineers & Health Care & 5.1 & $-$5.4\% & 25.6\% & $-$0.18 & 0.78 & $-$45.4\% & $-$8.6\% \\
K+S$^{\dagger}$ & Materials & 7.5 & $-$9.2\% & 42.5\% & $-$0.03 & 0.99 & $-$70.1\% & $-$24.7\% \\
TUI$^{\dagger}$ & Consumer & 2.7 & $-$9.3\% & 32.8\% & $-$0.25 & 1.09 & $-$50.5\% & $-$8.5\% \\
Commerzbank & Financials & 16.4 & $-$9.7\% & 46.2\% & $-$0.02 & 1.37 & $-$98.1\% & $-$24.0\% \\
Lanxess$^{\dagger}$ & Materials & 3.0 & $-$11.5\% & 27.5\% & $-$0.31 & 0.90 & $-$49.3\% & $-$7.0\% \\
Porsche SE$^{\dagger}$ & Financials & 4.8 & $-$15.6\% & 29.2\% & $-$0.51 & 1.05 & $-$62.1\% & $-$10.9\% \\
ProSiebenSat.1 Media$^{\dagger}$ & Telecom & 2.0 & $-$17.9\% & 26.3\% & $-$0.60 & 0.93 & $-$40.7\% & $-$14.5\% \\
Porsche AG$^{\dagger}$ & Consumer & 2.8 & $-$21.4\% & 29.3\% & $-$0.78 & 0.91 & $-$63.9\% & $-$7.2\% \\
Sartorius$^{\dagger}$ & Health Care & 4.1 & $-$21.5\% & 46.2\% & $-$0.34 & 1.15 & $-$70.7\% & $-$18.5\% \\
Zalando & Consumer & 5.1 & $-$23.6\% & 47.4\% & $-$0.38 & 1.37 & $-$82.0\% & $-$13.4\% \\
Scout24 & Technology & 1.1 & $-$24.9\% & 31.8\% & $-$0.81 & 0.48 & $-$40.7\% & $-$6.1\% \\
Puma$^{\dagger}$ & Consumer & 1.3 & $-$39.7\% & 39.8\% & $-$1.07 & 1.28 & $-$61.3\% & $-$9.0\% \\
Delivery Hero$^{\dagger}$ & Consumer & 1.8 & $-$41.7\% & 63.4\% & $-$0.52 & 1.23 & $-$82.5\% & $-$30.4\% \\
Hochtief & Industrials & 0.3 & $-$58.7\% & 33.8\% & $-$2.51 & 1.26 & $-$28.4\% & $-$9.4\% \\
HelloFresh$^{\dagger}$ & Consumer & 1.0 & $-$70.2\% & 66.6\% & $-$1.48 & 1.41 & $-$75.6\% & $-$14.0\% \\
\end{longtable}
}

\begin{thebibliography}{9}
\bibitem[Ang et al.(2006)]{ang2006} Ang, A., Hodrick, R., Xing, Y., Zhang, X. (2006). The cross-section of
volatility and expected returns. \emph{Journal of Finance} 61(1), 259--299.
\bibitem[Bailey and L\'opez de Prado(2014)]{bailey2014} Bailey, D., L\'opez de Prado, M. (2014). The
deflated Sharpe ratio. \emph{Journal of Portfolio Management} 40(5), 94--107.
\bibitem[Blitz and van Vliet(2007)]{blitz2007} Blitz, D., van Vliet, P. (2007). The volatility effect.
\emph{Journal of Portfolio Management} 34(1), 102--113.
\bibitem[Frazzini and Pedersen(2014)]{frazzini2014} Frazzini, A., Pedersen, L. (2014). Betting against
beta. \emph{Journal of Financial Economics} 111(1), 1--25.
\bibitem[George and Hwang(2004)]{george2004} George, T., Hwang, C. (2004). The 52-week high and momentum
investing. \emph{Journal of Finance} 59(5), 2145--2176.
\bibitem[Harvey et al.(2016)]{harvey2016} Harvey, C., Liu, Y., Zhu, H. (2016). \ldots and the cross-section
of expected returns. \emph{Review of Financial Studies} 29(1), 5--68.
\bibitem[Jegadeesh(1990)]{jegadeesh1990} Jegadeesh, N. (1990). Evidence of predictable behavior of
security returns. \emph{Journal of Finance} 45(3), 881--898.
\bibitem[Jegadeesh and Titman(1993)]{jegadeesh1993} Jegadeesh, N., Titman, S. (1993). Returns to buying
winners and selling losers. \emph{Journal of Finance} 48(1), 65--91.
\bibitem[Mertens(2002)]{mertens2002} Mertens, E. (2002). Comments on variance of the IID estimator in Lo
(2002). Working paper, University of Basel.
\end{thebibliography}
\end{document}
